\documentclass[book.tex]{subfiles}
\begin{document}
\chapter{Universal Ordinary Differential Equations and Residual Physics}

\label{chap:pinns}
\noindent\rule{11cm}{0.4pt} \\
\textbf{Prerequisites:}  \autoref{chap:pde}, \autoref{chap:numerical_pde}, \autoref{chap:neural_ode}  \\
\textbf{Difficulty Level:} ** \\
\noindent\rule{11cm}{0.4pt}
\vspace{5mm}

In previous chapters we have learned of a number of techniques to approximately solve differential equations using differentiable programs. Most such techniques directly approximate the solution $u$ to the differential equation by a suitable deep network. The universal ordinary differential equations \cite{bills2020universal}, \cite{rackauckas2020universal} approach instead seeks to use the deep network as a universal approximator that learns an error correction term that modifies a standard solution with a learned correction term.

%"""
%A central challenge in scientific machine learning is reconciling data that is at odds with simplified models without requiring "big data". In this work, universal differential equations (UDEs) are developed that augment scientific models with machine-learnable structures for scientifically-based learning. We show how UDEs can be utilized to discover previously unknown governing equations, accurately extrapolate beyond the original data, and accelerate model simulation, all in a time and data-efficient manner. As an application of UDEs, we use a machine-learning based physics-informed battery performance model to break the typically observed accuracy-computing cost trade-off. 
%"""

\section{Mathematical Framework} 

The fully general formulation of a universal differential equation is given by
\begin{align}
    \mathcal{N}[u(t), u(\alpha(t)), W(t), U_\theta(u, \beta(t))] = 0
\end{align}
where $\mathcal{N}$ is a possibly nonlinear operator, $u$, is a proposed solution, $\alpha(t)$ is a delay function, and $W(t)$ is the Wiener process.

Let $u_\theta(t)$ denote the universal differential equation. The loss function for such an equation is given by
\begin{align}
    \mathcal{L}(\theta) &= \sum_i \|u_{\theta}(t_i) - d_i\|
\end{align}
where $(t_1, d_1),\dotsc,(t_n,d_n)$ are given discrete datapoints. 


\section{Lotka-Volterra Equations}

As a concrete example, consider the Lotka-Volterra equations. We add on approximation terms $U_1, U_2$ to the two quantities being solved for in this equation.

\begin{align}
    \dot{x} &= \alpha x + U_1(\theta, x, y) \\
    \dot{y} &= - \theta_1 y + U_2(\theta, x, y)
\end{align}
\end{document}